% =============================================================================
% 1. ABSTRACT
% =============================================================================
\vspace{-2.5ex}

\begin{abstract}
Coding agents must integrate external tool returns into ongoing reasoning---a
capability that standard left-to-right pretraining on code exposes only in its
forward direction. We observe that the
\textit{action $\rightarrow$ observation $\rightarrow$ continuation} loop of a
coding agent is structurally isomorphic to a function call site, where a
caller binds arguments, a callee returns a value computed elsewhere, and
downstream code consumes that value. This conditioning structure exists at
internet scale in ordinary code. We exploit it through
\emph{function-aware fill-in-the-middle (FIM) mid-training}: a self-supervised
objective that masks functions selected via program dependency graph
analysis and a complexity--inferability double criterion. We mid-train
Qwen2.5-Coder-Instruct (7B/14B) and Qwen3-8B on a 2.6B-token decontaminated
corpus drawn from 968 GitHub repositories, then apply existing agentic
post-training pipelines.
Mid-training improves SWE-Bench-Verified by $+2.8/+3.0$ at 7B/14B and by
$+3.2$ on Qwen3-8B; SWE-Bench-Lite gains are $+3.7/+4.0/+5.4$ on the same
models. The improvement holds across two post-training pipelines (R2E-Gym,
SWE-Smith) and on a non-Qwen2.5 base (Qwen3-8B with SWE-Lego). Beyond
in-domain gains, mid-training also mitigates the capability erosion that
agentic post-training otherwise inflicts on non-agent coding (e.g.,
LiveCodeBench) and non-coding tool-use benchmarks ($\tau$-bench, BFCL):
although the mid-training corpus contains Python code only, the
function-call inductive bias survives post-training and yields consistent gains.
\end{abstract}

\vspace{-2ex}